\setcounter{section}{27}
\section{Sistemas de la Información}

Nombre completo: Definición, estructura, y dimensionamiento eficiente de los sistemas de información.

\localtableofcontents

\subsection{Sistema de información}
Es un conjunto de elementos orientados al tratamiento y administración de datos e informaciónm, organizados y listos para su uso psterior, generados para cubrir una necesidad u objetivo. Los elementos pueden ser: Personas, Actividades, Datos, recursos materiales. Todos estos procesos interactúan para procesar los datos (incluidos procesos manuales y automáticos) y dan lugar a información mas elaborada. 

El término se usa erróneamente como sistema de información \textbf{informático}, en parte porque la mayoría de los recursos son sistemas informáticos. Pero puede no disponer de estos recursos.

Se entiende un sistema de información como \textit{el conjunto de tecnologías, procesos aplicaciones de negocio y software disponibles para las personas dentro de una organización}.

\subsubsection{Otros sistemas de información}

En diferentes campos, un sistema de información puede hacer referencia a otros conceptos.

\begin{description}
    \item[Geografía y Cartografía] se usa un \textbf{Sistema de Información Geográfica} para tratar con información georreferenciada.
    \item[Informática] cualquier sistema informático que se use para tratar datos.
    \item[Matemáticas] son sistemas de atributo-valor
    \item[Representación del conocimiento] Consta de tres componentes, humano, tecnológico y organizacional y se define en términos de tres niveles de semiótica:
    \begin{description}
        \item[Sintáxis] datos que se pueden procesar automáticamente
        \item[Semántico] Datos procesados por un individuo.
        \item[Pragmático] Cuando un individuo conoce, entiende y evalúa la información
    \end{description}
    \item[Seguridad computacional] Debe tener estos tres componentes
    \begin{description}
        \item[Estructura] Interfaces de intercambio y repositorios de datos
        \item[Canales] Que conectan repositorios entre sí
        \item[Comportamiento] Mensajes que acarréan contenido/significado y Servicios
    \end{description}
\end{description}



\subsection{Componentes básicos}

Por tanto un sistema de información debe cumplic con los siguientes componentes interactuando entre sí:
\begin{description}
    \item[Hardware] El equípo físico utilizado para procesar y almacenar datos
    \item[Software] y procedimientos utilizados para transformar y extraer información
    \item[Datos] que representan las actividades de la empresa
    \item[Red] que permite compartir recursos entre los distintos dispositivos
    \item[Personas] que desarrollan, mantienen y utilizan el sistema
\end{description}

Los sistemas de información son una combinación de tres partes: las personas, los procesos y los equipos.

\subsection{Actividades}

Entrada, almacenamiento, transformación y salida. A veces la salida se conecta con otra (o la misma) entrada. De esta forma se conectan varios sistemas de información entre sí, generando un sistema de información (mas grande).

Los sistemas de información son evaluados en base a su utilidad, es decir, lo que facilita o mejora el desempeño de un individuo.

Algunos parámetros son: relevancia, comprensibilidad, completud, tiempo de respuesta, facilidad de uso, confiabilidad y flexibilidad.

\subsection{Ciclo de vida}

\begin{description} %El original está ordenado alfabétimcamente, a ver si lo sacamos
    % Esto es un waterfall como una casa, agrupamos en esos pasos

    % Requisitos
    \item[Conocimiento de la organización]
    \item[Identificación de problemas y oportunidades] 
    \item[Determinar necesidades] 
    \item[Diagnóstico]
    \item[Propuesta] 
    % Diseño
    \item[Diseño de sistema] 
    % Implementación & pruebas
    \item[Codificación]
    % Despliegue
    \item[Implementación] 
    % Mantenimiento
    \item[Mantenimiento] 
\end{description}

\subsection{Aplicación}

El producto actual de muchas empresas es el conocimiento mas que el objeto. El énfasis ha cambiado hacia los servicios. Siendo la información, uno de los mayores activos de una empresa. 